% !TeX root = ../mathematical_methods.tex

\section{Introduzione}
        E' opportuno introdurre questa nuova sezione con una propria definizione di \textit{segnale}:

        \begin{definition}[Definizione di segnale]
            Un segnale è una qualsiasi funzione di variabili indipendenti, ovvero
            \[
                f: \Omega \subseteq \mathbb{R}^{n} \rightarrow \mathbb{R}
            \]
            \[
                \mathbf{t} \in \Omega \mapsto f(\mathbf{t})
            \]
        \end{definition}

        Da un punto di vista fisico, un segnale è \textit{una variazione di una precisa quantità fisica in funzione di tempo, spazio, o qualsiasi altra variabile indipendente}.
        Noi lavoreremo per la maggior parte con segnali tempo-varianti.

        Determinato il significato, matematico e fisico, del segnale, possiamo procedere a definire una serie di elementi utili nell'analisi di quest'ultimi.

        \begin{definition}[Definizione di supporto]
            Sia $f: \Omega \subseteq \mathbb{R} \rightarrow \mathbb{R}$ un segnale.
            Si dice supporto di $f$ il seguente insieme
            \[
                S(f) = \overline{\{t \in \Omega \; : \; f(t) \neq 0\}}
            \]
        \end{definition}

        In altre parole, il supporto del segnale è la chiusura (nel senso topologico) dell'insieme degli elementi in cui il segnale non è nullo.

        \begin{definition}[Definizione di segnale localmente sommabile]
            Sia $f: \Omega \subseteq \mathbb{R} \rightarrow \mathbb{R}$ un segnale.
            Questi si dice localmente sommabile se $\forall \: a < b < \infty \in \Omega$, si ha che
            \[
                \int_{a}^{b} \abs{f(t)} \: dt < \infty
            \]
        \end{definition}

        \begin{definition}[Definizione di segnale sommabile]
            Sia $f: \Omega \subseteq \mathbb{R} \rightarrow \mathbb{R}$ un segnale.
            Questi si dice sommabile, o assolutamente integrabile, se
            \[
                \int_{\Omega} \abs{f(t)} \: dt < \infty
            \]
        \end{definition}


    \section{Teoria delle distribuzioni}
        Denotiamo con il termine distribuzioni tutti quelli oggetti matematici che si \textit{comportano} come funzioni, ovvero oggetti che mappano un input ad un determinato output.
        In altre parole, le distribuzioni rappresentano una \textit{generalizzazione} del concetto di funzione.
        Questi oggetti sono caratterizzati, ovvero sono determinati nella loro interezza, a seconda dell'effetto che hanno, intesi come funzionali lineari, su uno spazio di appositamente scelte funzioni lisce su $\mathbb{R}$, il cui supporto è compatto, chiamate \textit{funzioni di test}.

        \begin{definition}[Definizione di spazio di funzioni di test]
            Sia $\mathcal{D} = C^{\infty}_{c}(\mathbb{R})$, dove per ogni $\phi \in \mathcal{D}$, si ha che
            \begin{itemize}
                \item $\phi$ è infinitamente differenziabile.
                \item $S(\phi)$ è compatto (per il teorema di Heine-Borel, in $\mathbb{R}^{n}$, questo equivale a dire che $S(\phi)$ è chiuso e limitato).
            \end{itemize}

            Allora $\mathcal{D}$ è uno spazio vettoriale su $\mathbb{R}$.
        \end{definition}

        \begin{definition}[Definizione di distribuzione]
            Una distribuzione $T: \mathcal{D} \rightarrow \mathbb{R}$ è, formalmente, un funzionale lineare continuo (nel senso topologico del termine) definito su uno spazio $\mathcal{D}$ di funzioni di test.
            In altre parole, essa è un elemento dello spazio duale topologico $\mathcal{D}^{'}$.
        \end{definition}


        \begin{proposition}
            Sia $T \in \mathcal{D}^{'}$ una distribuzione.
            Si definisce $T = 0$ per un insieme $U \subseteq \mathbb{R}$ aperto se $\forall \phi \in \mathcal{D}$ tale per cui $S(\phi) \subseteq U$, si ha che $\langle T, \phi \rangle = 0$.
            Allora, data una arbitraria unione di insiemi $V = \bigcup \{U_{\alpha}\}_{\alpha \in A}$ tali per cui $T = 0$ per ogni $U_{\alpha}$, si ha che $T = 0$ per $V$.
        \end{proposition}

        \begin{proof}[Dimostrazione]
            Si consideri una funzione di test $\phi \in \mathcal{D}$ tale per cui $S(\phi) \subseteq V$.
            Dato che $S(\phi)$ è un insieme compatto, per Heine-Borel, si ha che, per ogni copertura aperta di $S(\phi)$, questa ammette sottocopertura finita.
            Di conseguenza, è possibile determinare $k$ elementi di $\{U_{\alpha}\}_{\alpha \in A}$, tali per cui
            \[
                S(\phi) \subseteq \bigcup_{i = 1}^{k} U_{i}
            \]
            Per ognuno di questi insiemi $U_{i}$ è possibile determinare una funzione di test $\psi_{i}$ tale per cui $S(\psi_{i}) \subseteq U_{i}$ e che $\forall x \in S(\phi), \; \exists \psi_{i} \; | \; \psi_{i}(x) > 0$.
            Si consideri allora una funzione $\psi_{0} \in C^{\infty}(\mathbb{R})$ tale per cui $\psi_{0}(x) = 0$ per tutti gli $x \in S(\phi)$ e $\psi_{0}(x) > 0$ per tutti gli $x \notin S(\phi)$.
            Allora
            \[
                \kappa(x) = \sum_{i = 0}^{k} \psi_{i}(x) \neq 0, \; \forall x \in \mathbb{R} \implies \frac{1}{\kappa} \in C^{\infty}(\mathbb{R})
            \]
            Si può osservare quanto segue
            \[
                \rho_{i}(x) = \frac{\psi_{i}(x)}{\kappa(x)} \in C^{\infty}_{c}(U_{i})
            \]
            Se $x \in S(\phi)$, si ha che
            \[
                \sum_{i = 1}^{k} \rho_{i} (x) = \sum_{i = 1}^{k} \frac{\psi_{i}(x)}{\kappa(x)} = \sum_{i = 1}^{k} \frac{\psi_{i}(x)}{\sum_{i = 1}^{k} \psi_{i}(x) + 0} = \frac{\sum_{i = 1}^{k} \psi_{i}(x)}{\sum_{i = 1}^{k} \psi_{i}(x)} = 1
            \]
            $\sum_{i = 1}^{k} \rho_{i}(x)$ è una combinazione lineare di funzioni test aventi supporto, rispettivamente, negli insiemi $U_{i}$.
            Allora
            \[
                \phi = \phi \cdot 1 = \phi \cdot \sum_{i = 1}^{k} \rho_{i} = \sum_{i = 1}^{k} \phi \cdot \rho_{i} \implies \langle T, \phi \rangle = \langle T, \sum_{i = 1}^{k} \phi \cdot \rho_{i} \rangle = \sum_{i = 1}^{k} \langle T, \phi \cdot \rho_{i} \rangle
            \]
            Dato che $\phi \cdot \rho_{i} \in \mathcal{D}$ e $S(\phi \cdot \rho_{i}) \subseteq U_{i}$
            \[
                \sum_{i = 1}^{k} \langle T, \phi \cdot \rho_{i} \rangle = \sum_{i = 1}^{k} 0 = 0
            \]
        \end{proof}

        \begin{definition}[Definizione di supporto di una distribuzione]
            Sia $T \in \mathcal{D}^{'}$ una distribuzione.
            Sia allora $V = \bigcup \{U \; : \; U \; \textnormal{aperto tale per cui} \; T = 0\}$.
            Il supporto di $T$ diventa
            \[
                S(T) = \mathbb{R} \setminus V
            \]
        \end{definition}

        Le distribuzioni così definite ci conferiscono una serie di proprietà particolarmente comode:

        \begin{itemize}
            \item Linearità secondo distribuzione: $\langle \alpha T + \beta S, \phi\rangle = \alpha \langle T, \phi \rangle + \beta \langle S, \phi \rangle$
            \item Linearità secondo funzione di test: $\langle T, \alpha \phi + \beta \psi \rangle = \alpha \langle T, \phi \rangle + \beta \langle T, \psi \rangle$
        \end{itemize}

        E' importante sottolineare che funzioni localmente integrabili su $\mathbb{R}$ sono a loro volta distribuzioni mediante l'operazione di integrazione.
        In altre parole, data una funzione localmente $f$ integrabile su $\mathbb{R}$, è possibile definire la seguente distribuzione $F$
        \[
            F \in \mathcal{D}^{*}, \; F(\phi) := \int_{\mathbb{R}} f(t) \phi(t) \; dt
        \]

        In virtù di ciò, definiamo le seguenti operazioni:

        \begin{definition}[Traslazione di distribuzione]
            Sia $T \in \mathcal{D}^{*}$ una distribuzione.
            Definiamo la traslazione di $T$ di un $\tau \in \mathbb{R}$ come una distribuzione $T_{\tau}$ tale per cui
            \[
                \langle T_{\tau}, \phi \rangle = \langle T, \phi(t - \tau) \rangle
            \]
        \end{definition}

        \begin{definition}[Scala di distribuzione]
            Sia $T \in \mathcal{D}^{*}$ una distribuzione.
            Definiamo la scala di $T$ secondo un $\alpha > 0$ come una distribuzione $T(\alpha t)$ tale per cui
            \[
                \langle T(\alpha t), \phi \rangle = \frac{1}{\alpha} \langle T, \phi(\alpha^{-1} t)\rangle
            \]
        \end{definition}

        \begin{definition}{Derivata di distribuzione}
            Sia $T \in \mathcal{D}^{*}$ una distribuzione.
            Definiamo la derivata di distribuzione $T^{(n)}$ come
            \[
                \langle T^{(n)}, \phi \rangle = (-1)^{n} \langle T, \phi^{(n)} \rangle
            \]
        \end{definition}

        Grazie alla definizione di derivata di distribuzione, possiamo verificare che questo operatore rispetti la regola di Leibnitz.
        Sia $T \in \mathcal{D}^{*}$ una distribuzione e sia $f \in C^{\infty}(\mathbb{R})$, allora
        \[
            \langle (T f)', \phi \rangle = -\langle T f, \phi' \rangle = \langle T, -\phi' f \rangle
        \]
        \[
            \langle T' f + T f', \phi \rangle = \langle T', \phi f \rangle + \langle T, \phi f' \rangle = \langle T, \phi f' - (\phi f)' \rangle = \langle T, \phi f' - \phi f' - \phi' f \rangle = \langle T, - \phi' f \rangle = \langle (T f)', \phi \rangle
        \]



    \section{Operazioni sulle distribuzioni}
        \begin{definition}[Definizione di convoluzione tra distribuzione e funzione test]
            Sia $T \in \mathcal{D}^{*}$ una distribuzione e sia $\phi \in \mathcal{D}$ una funzione di test.
            Si definisce allora la convoluzione $T * \phi$ come segue
            \[
                (T * \phi) (x) = \langle T, \tau_{x} \overline{\phi} \rangle
            \]

            Se $T$ ha supporto compatto, allora anche la risultante convoluzione ha supporto compatto, ed è quindi una funzione di test, ovvero $T * \phi \in \mathcal{D}$.
        \end{definition}

        \begin{definition}[Definizione di convoluzione fra distribuzioni]
            Siano $T, S \in \mathcal{D}^{*}$ due distribuzioni tale per cui $S$ ha supporto compatto.
            Si definisce allora la convoluzione $T * S \in \mathcal{D}^{*}$ come segue
            \[
                \langle T * S, \phi \rangle = \langle T, \overline{S} * \phi \rangle = \langle T, \langle \overline{S}, \tau_{x} \overline{\phi} \rangle \rangle = \langle T, \langle S, \tau_{x} \phi \rangle \rangle = \langle T(x), \langle S(y), \phi (x + y) \rangle \rangle
            \]
        \end{definition}

        Qualora dovessimo trattare di distribuzioni determinate da segnali regolari, attingiamo alla seguente definizione integrale di convoluzione

        Siano $f, g: \mathbb{R} \rightarrow \mathbb{R}$ due segnali regolari.
        Si definisce allora la convoluzione $f * g: \mathbb{R} \rightarrow \mathbb{R}$ fra $f, g$ come
        \[
            (f * g) (t) = \int_{-\infty}^{\infty} f(\tau)g(t - \tau) \: d\tau \iff \int_{-\infty}^{\infty} \abs{f(\tau)} \abs{g(t - \tau)} \: d\tau < \infty, \; \forall t \in \mathbb{R}
        \]

        Alcune delle proprietà della convoluzione sono
        \begin{enumerate}
            \item Commutatività:
            \[
                f * g = \int_{-\infty}^{\infty} f(\tau) g(t - \tau) \: d\tau = - \int_{\infty}^{-\infty} f(t - u) g(u) \: du = \int_{-\infty}^{\infty} g(u) f(t - u) \: du = g * f
            \]

            \item Bilinearità:
            \[
                (\alpha f + \beta g) * (\lambda h + \mu k) = \alpha \lambda (f * h) + \alpha \mu (f * k) + \beta \lambda (g * h) + \beta \mu (g * k)
            \]

            \item Neutralità secondo delta di Dirac:
            \[
                f * \delta  = f
            \]
        \end{enumerate}


        Per quanto riguarda i sistemi lineari tempo-invarianti (LTI), la convoluzione rappresenta lo strumento attraverso il quale esprimiamo l'uscita rispetto ad un segnale in ingresso
        \[
            x(t) \mapsto y(t) = x * h = \int_{-\infty}^{\infty} x(\tau) h(t - \tau) \: dt
        \]
        dove h rappresenta la \textit{risposta all'impulso} del sistema LTI, ovvero la risposta ad una perturbazione assimilabile ad una delta di Dirac.

        Dimostriamo questo fatto nella seguente maniera.
        Dinanzi tutto, partiamo dal pressupposto che, dato un segnare regolare ad energia finita $x(t)$, si ha che
        \[
            x * \delta = x \implies x(t) = \int_{-\infty}^{\infty} x(\tau) \delta(t - \tau) \: d\tau
        \]

        Possiamo allora rappresentare la risposta del sistema ad un generico segnale in ingresso come un operatore lineare $\mathcal{H}$.
        Di conseguenza
        \[
            y(t) = \mathcal{H}(x(t)) = \mathcal{H}\left(\int_{-\infty}^{\infty} x(\tau) \delta(t - \tau) \: d\tau\right)
        \]
        per la linearità dell'operatore $\mathcal{H}$
        \[
            \mathcal{H}\left(\int_{-\infty}^{\infty} x(\tau) \delta(t - \tau) \: d\tau\right) = \int_{-\infty}^{\infty} \mathcal{H}(x(\tau) \delta(t - \tau)) \: d\tau = \int_{-\infty}^{\infty} x(\tau) \mathcal{H}(\delta(t - \tau)) \: d\tau
        \]
        denotando la risposta all'impulso di Dirac del sistema come $h(t) = \mathcal{H}(\delta(t))$, si arriva alla conclusione che
        \[
            y(t) = \int_{-\infty}^{\infty} x(\tau) h(t - \tau) \: d\tau
        \]

        \begin{definition}[Definizione di sistema causale]
            Sia $h(t): \mathbb{R} \rightarrow \mathbb{R}$ la risposta all'impulso di un sistema LTI.
            Il sistema si dice causale se il supporto di $h$ è contenuto in $[0, +\infty)$ e, dato un segnale regolare $x(t)$ in ingresso, vale la seguente relazione
            \[
                \int_{-\infty}^{\infty} x(\tau) h(t - \tau) \: d\tau = \int_{-\infty}^{t} x(\tau) h(t - \tau) \: d\tau
            \]
        \end{definition}



    \section{Trasformata di Fourier}
        Sia $x(t) \in L^{1}(\mathbb{R})$ un segnale.
        Si definisce la sua trasformata di Fourier nel seguente modo
        \[
            \hat{x}(\omega) = \int_{-\infty}^{\infty} x(t)e^{-i \omega t} \: dt
        \]

        Dimostriamo allora come è possibile ottenere $\hat{x}$ a partire da una serie di Fourier.
        Si immagini di considerare $x(t)$ come fosse un segnale periodico a periodo $\mathbb{R}$.
        Di conseguenza, questi ha una serie di Fourier i cui coefficienti sono dati da
        \[
            c_{k} = \frac{1}{T} \int_{-T/2}^{T/2} x(t) e^{-i k \omega_{0} t} \: dt \implies T c_{k} \int_{-T/2}^{T/2} x(t) e^{-i k \omega_{0} t} \: dt
        \]

        Prendendo il limite per $T \rightarrow \infty$, si ha che
        \[
            \lim_{T \rightarrow \infty} (T c_{k}) = \lim_{T \rightarrow \infty} \; \int_{-T/2}^{T/2} x(t) e^{-i k \omega_{0} t} \: dt = \int_{-\infty}^{\infty} x(t) e^{-i \omega t} \: dt = \hat{x}(\omega)
        \]

        Alcune delle proprietà principali della trasformata di Fourier sono:
        \begin{enumerate}
            \item Linearità:
            \[
                \mathcal{F}\{x + h\} = \int_{-\infty}^{\infty} (x(t) + h(t)) e^{-i \omega t} \: dt = \int_{-\infty}^{\infty} x(t) e^{-i \omega t} \: dt + \int_{-\infty}^{\infty} h(t) e^{-i \omega t} \: dt = \mathcal{F}\{x\} + \mathcal{F}\{h\}
            \]
            \item Modulazione del segnale in ingresso:
            \[
                \mathcal{F}(x(t) e^{i \omega_{0} t}) = \int_{-\infty}^{\infty} x(t) e^{i (\omega_{0} - \omega) t} \: dt = \mathcal{F}\{x\}(\omega - \omega_{0})
            \]
            \item Traslazione del segnale in ingresso:
            \[
                \mathcal{F}\{x(t - t_{0})\} = \int_{-\infty}^{\infty} x(t - t_{0}) e^{-i \omega t} \: dt = \int_{-\infty}^{\infty} x(u) e^{-i \omega (u + t_{0})} \: dt = e^{-i \omega t_{0}} \mathcal{F}\{x(t)\}
            \]
            \item Riscalamento del segnale in ingresso:
            \[
                \forall \alpha \neq 0, \; \mathcal{F}\{x(\alpha t)\} = \int_{-\infty}^{\infty} x(\alpha t) e^{-i \omega t} \: dt = \frac{1}{\alpha} \int_{-\infty}^{\infty} x(u) e^{-i \omega \frac{u}{\alpha}} \: dt = \frac{1}{\alpha} \mathcal{F}\{x\}\left(\frac{\omega}{\alpha}\right)
            \]

            \item Derivata del segnale in ingresso: Sia $x(t) \in L^{1}(\mathbb{R})$ tale per cui $x'(t) \in L^{1}(\mathbb{R})$
            \[
                \mathcal{F}\{x'(t)\} = \int_{-\infty}^{\infty} x'(t) e^{-i \omega t} \: dt = x(t) e^{-i \omega t} \rvert_{-\infty}^{\infty} - (-i \omega) \int_{-\infty}^{\infty} x(t) e^{-i \omega t} \: dt
            \]
            In virtù del fatto che $x(t)$ abbia derivata in $L^{1}(\mathbb{R})$, si ha che per $t \rightarrow \pm \infty$, $x(t) \rightarrow 0$, quindi
            \[
                \mathcal{F}\{x'(t)\} = 0 + i \omega \int_{-\infty}^{\infty} x(t) e^{-i \omega t} \: dt = i \omega \mathcal{F}\{x(t)\}
            \]

            \item Derivata di trasformata:
            \[
                \mathcal{F}'\{x(t)\}(\omega) = \frac{d}{d\omega} \int_{-\infty}^{\infty} x(t) e^{-i \omega t} \: dt = \int_{-\infty}^{\infty} x(t) \frac{d}{d\omega} e^{-i \omega t} \: dt = \int_{-\infty}^{\infty} x(t) \cdot -it e^{-i \omega t} \: dt
            \]
            \[
                = -i \int_{-\infty}^{\infty} (t \cdot x(t)) e^{-i \omega t} \: dt = -i \mathcal{F}\{t x(t)\}
            \]

            \item Convoluzione di segnali in ingresso:
            \[
                \mathcal{F}\{x(t) * y(t)\} = \int_{-\infty}^{\infty} \left(\int_{-\infty}^{\infty} x(\tau) y(t - \tau) \: d\tau\right) e^{-i \omega t} \: dt = \int_{-\infty}^{\infty} x(\tau) \left[\int_{-\infty}^{\infty} y(t - \tau) e^{-i \omega t} \: dt\right] \: d\tau =
            \]
            \[
                \int_{-\infty}^{\infty} x(\tau) \left[\int_{-\infty}^{\infty} y(u) e^{-i \omega (u + \tau)} \: du\right] \: d\tau = \int_{-\infty}^{\infty} x(\tau) e^{-i \omega \tau} \: d\tau \int_{-\infty}^{\infty} y(u) e^{-i \omega u} \: du = \mathcal{F}\{x(t)\} \mathcal{F}\{y(t)\}
            \]
        \end{enumerate}

        Un'altra peculiare proprietà della trasformata di Fourier è la seguente, ovvero
        \[
            \mathcal{F}\{x(-t)\}(\omega) = \mathcal{F}\{x(t)\}(-\omega)
        \]
        La cui dimostrazione è suggerita dal fatto che
        \[
            \mathcal{F}\{x(-t)\} = \int_{-\infty}^{\infty} x(-t) e^{-i \omega t} \: dt = \int_{-\infty}^{\infty} x(u) e^{-i (-\omega) u} \: du = \mathcal{F}\{x(u)\}(-\omega)
        \]

        \begin{lemma}[Lemma di Riemann-Lebesgue]
            Sia $x(t) \in L^{1}(\mathbb{R})$ un segnale assolutamente integrabile tale per cui $x'(t) \in L^{1}(\mathbb{R})$.
            Allora
            \[
                \lim_{\omega \rightarrow \pm \infty} \hat{x}(\omega) = 0
            \]
        \end{lemma}

        \begin{proof}[Dimostrazione]
            Si consideri l'espressione intera della trasformata
            \[
                \hat{x}(\omega) = \int_{-\infty}^{\infty} x(t) e^{-i \omega t} \: dt \implies \abs{\int_{-\infty}^{\infty} x(t) e^{-i \omega t} \: dt} \leq \int_{-\infty}^{\infty} |x(t)| |e^{-i \omega t}| \: dt \leq \int_{-\infty}^{\infty} |x(t)| \: dt < \infty, \; \forall \omega \in \mathbb{R}
            \]

            Applicando l'integrazione per parti e opportunamente scegliendo $dv = e^{-i \omega t} \: dt$, $u = x(t)$, si ottiene
            \[
                \int_{-\infty}^{\infty} x(t) e^{-i \omega t} \: dt = x(t) e^{-i \omega t} \left. \frac{1}{-i \omega} \right\rvert_{-\infty}^{\infty} + \frac{1}{i \omega} \int_{-\infty}^{\infty} x'(t) e^{-i \omega t} \: dt
            \]
            \[
                = \frac{1}{i \omega} \left[-x(t) e^{-i \omega t} \rvert_{-\infty}^{\infty} + \int_{-\infty}^{\infty} x'(t) e^{-i \omega t} \: dt\right]
            \]

            Osservando attentamente il comportamento di $x'(t)$, si ha
            \[
                x'(t) \in L^{1}(\mathbb{R}) \implies -x(t) e^{-i \omega t} \rvert_{-\infty}^{\infty} = 0, \forall \omega \in \mathbb{R}
            \]

            Per quanto riguarda il termine integrale, invece
            \[
                \int_{-\infty}^{\infty} x'(t) e^{-i \omega t} \: dt = \mathcal{F}\{x'(t)\}(\omega) \implies \abs{\mathcal{F}\{x'(t)\}(\omega)} < \infty, \; \forall \omega \in \mathbb{R}
            \]

            La somma diventa allora
            \[
                -x(t) e^{-i \omega t} \rvert_{-\infty}^{\infty} + \int_{-\infty}^{\infty} x'(t) e^{-i \omega t} \: dt = 0 + z, \; |z| < M, \; \forall \omega \in \mathbb{R}
            \]

            per qualche $M \in \mathbb{R}$. Mettendo il tutto assieme, si ottiene
            \[
                \lim_{\omega \rightarrow \pm \infty} \hat{x}(\omega) = \lim_{\omega \rightarrow \pm \infty} \frac{1}{i \omega} \cdot z = 0
            \]
        \end{proof}

        \subsection{Antitrasformata di Fourier}
            \begin{theorem}[Esistenza dell'antitraformata]
                Sia $x(t) \in L^{1}(\mathbb{R})$ continua, e sia la sua traformata $\mathcal{F}\{x(t)\}$ a sua volta in $L^{1}(\mathbb{R})$.
                Allora $\mathcal{F}\{x(t)\}$ è antitrasformabile con antitrasformata
                \[
                    x(t) = \frac{1}{2 \pi} \int_{-\infty}^{\infty} \mathcal{F}\{x(t)\} e^{i \omega t} \: d\omega = \frac{1}{2 \pi} \mathcal{F}\{\mathcal{F}\{x(t)\}\}(-t)
                \]
            \end{theorem}

        \subsection{Estensione della trasformata di Fourier alle distribuzioni}
            Sia $x(t) \in L^{1}(\mathbb{R})$, e sia $\hat{x}(\omega)$ la sua trasformata.
            Si ricordi che, conseguenza del fatto che $x(t)$ sia assolutamente integrabile si ottiene
            \[
                \hat{x}(\omega) = \int_{-\infty}^{\infty} x(t) e^{-i \omega t} \: dt \implies \abs{\int_{-\infty}^{\infty} x(t) e^{-i \omega t} \: dt} \leq \int_{-\infty}^{\infty} |x(t)| |e^{-i \omega t}| \: dt \leq \int_{-\infty}^{\infty} |x(t)| \: dt < \infty, \; \forall \omega \in \mathbb{R}
            \]

            Allora, data una $\phi \in \mathcal{D}$
            \[
                \abs{\phi(\omega) \hat{x}(\omega)} \leq k \abs{\phi(\omega)}
            \]

            Dato che la funzione $\phi(\omega) \hat{x}(\omega)$ è limitata e possiede un supporto compatto, essa è a sua volta assolutamente integrabile ($\phi(\omega) \hat{x}(\omega) \in L^{1}(\mathbb{R})$).
            Adesso, è necessario introdurre un nuovo spazio di funzioni:

            \begin{definition}[Spazio di Schwartz]
                Lo spazio di Schartz (anche chiamato lo spazio delle funzioni "rapidamente decrescenti") è lo spazio complesso $\mathcal{S}$ di tutte le funzioni complesse di classe $C^{\infty}(\mathbb{R})$ tali per cui
                \[
                    \psi \in \mathcal{S} \implies \lim_{t \rightarrow \pm \infty} \psi^{(n)}(t)Q(t) = 0, \; \forall Q \in \mathbb{R}[t], \; \forall n \in \mathbb{N}
                \]
            \end{definition}

            Armati di questa nuova definizione, prendiamo allora in considerazione il funzionale lineare $\langle \hat{x}, \cdot \rangle \in \mathcal{S}^{*}$
            \[
                \langle \hat{x}, \phi \rangle = \int_{-\infty}^{\infty} \int_{-\infty}^{\infty} x(t) e^{-i \omega t} \: dt \: \phi(\omega) \: d\omega = \int_{-\infty}^{\infty} x(t) \int_{-\infty}^{\infty} e^{-i \omega t} \phi(\omega) \: d\omega \: dt = \int_{-\infty}^{\infty} x(t) \hat{\phi}(t) \: dt = \langle x, \hat{\phi} \rangle
            \]

            Questo ci permette di definire la trasformata di Fourier per le distribuzioni, nello specifico distribuzioni temperate, ovvero i funzionali lineari nello spazio di Schwartz ($\mathcal{S}^{*}$).

            Inoltre, è possibile dimostrare la formula di inversione anche per la trasformata di distribuzione
            \[
                \frac{1}{2 \pi} \langle \hat{\hat{x}}(-t), \phi \rangle = \frac{1}{2 \pi} \langle \hat{x}, \hat{\phi}(-t) \rangle = \langle x, \frac{1}{2 \pi} \hat{\hat{\phi}}(-t) \rangle = \langle x, \phi \rangle
            \]
            Di conseguenza, seguono, nelle trasformate di distribuzioni, tutte le proprietà delle trasformate di Fourier su funzioni assolutamente integrabili.

            Alcune trasformate di distribuzioni note sono:
            \begin{enumerate}
                \item $\displaystyle \mathcal{F}\{u(t)\} = \PV \left\{\frac{1}{i \omega}\right\} + \pi \delta(\omega)$
                \item $\displaystyle \mathcal{F}\{\delta(t)\} = 1$
            \end{enumerate}
